\section{Lubricantes}

\subsection{Presentación del caso}
Una empresa industrial fabrica tres tipos de aceites lubricantes: aceite de motor, aceite hidráulico y aceite de transmisión. 
Las contribuciones unitarias al beneficio son, respectivamente, 4, 6 y 5 unidades monetarias por m$^3$ producido.  

La producción se organiza en tres procesos principales: calentamiento de materias primas, mezclado y envasado, que requieren la dedicación de 1, 2 y 4 trabajadores respectivamente.  

Cada tipo de aceite emplea los siguientes tiempos de proceso (en minutos por m$^3$):

\begin{center}
\begin{tabular}{lccc}
\toprule
 & Calentamiento & Mezclado & Envasado \\
\midrule
Aceite de motor & 4 & 6 & 8 \\
Aceite hidráulico & 4 & 10 & 15 \\
Aceite de transmisión & 4 & 12 & 5 \\
\bottomrule
\end{tabular}
\end{center}

Los tiempos máximos disponibles al día son 240 minutos de calentamiento, 480 minutos de mezclado y 960 minutos de envasado.  

Se pide plantear un modelo de programación lineal que maximice el beneficio.

\subsection{Formulación explícita}

\paragraph{Variables}\mbox{}\\
\begin{tabular}{lcl}
$x_m$ & : & cantidad (m$^3$) de aceite de motor producido \\
$x_h$ & : & cantidad (m$^3$) de aceite hidráulico producido \\
$x_t$ & : & cantidad (m$^3$) de aceite de transmisión producido \\
\end{tabular}

\paragraph{Restricciones}\mbox{}\\

Tiempo máximo de calentamiento:
\begin{equation}
4x_m + 4x_h + 4x_t \leq 240
\end{equation}

Tiempo máximo de mezclado:
\begin{equation}
6x_m + 10x_h + 12x_t \leq 480
\end{equation}

Tiempo máximo de envasado:
\begin{equation}
8x_m + 15x_h + 5x_t \leq 960
\end{equation}

\paragraph{Función objetivo}\mbox{}\\
\begin{equation}
\mbox{max.}\; z = 4x_m + 6x_h + 5x_t
\end{equation}

\paragraph{Modelo completo}\mbox{}\\
\begin{align}
\mbox{max.}\quad & 4x_m + 6x_h + 5x_t \notag \\
\mbox{s.a.}\quad 
& 4x_m + 4x_h + 4x_t \leq 240 \notag \\
& 6x_m + 10x_h + 12x_t \leq 480 \notag \\
& 8x_m + 15x_h + 5x_t \leq 960 \notag \\
& x_m, x_h, x_t \geq 0
\end{align}

\subsection{Formulación generalizada}

\paragraph{Conjuntos}\mbox{}\\
\begin{tabular}{lcl}
$\mathcal{A}$ & : & tipos de aceites lubricantes \\
$\mathcal{P}$ & : & procesos de producción (calentamiento, mezclado, envasado) \\
\end{tabular}

\paragraph{Parámetros}\mbox{}\\
\begin{tabular}{lcl}
$B_a$ & : & beneficio unitario por m$^3$ de aceite $a \in \mathcal{A}$ \\
$T_{pa}$ & : & tiempo requerido (min) del proceso $p \in \mathcal{P}$ por cada m$^3$ de aceite $a \in \mathcal{A}$ \\
$D_p$ & : & tiempo disponible (min) para el proceso $p \in \mathcal{P}$ \\
\end{tabular}

\paragraph{Variables}\mbox{}\\
\begin{tabular}{lcl}
$x_a$ & : & cantidad de aceite $a \in \mathcal{A}$ a producir (m$^3$) \\
\end{tabular}

\paragraph{Restricciones}\mbox{}\\
Para cada proceso $p \in \mathcal{P}$:
\begin{equation}
\sum_{a \in \mathcal{A}} T_{pa}\,x_a \leq D_p
\end{equation}

\paragraph{Función objetivo}\mbox{}\\
\begin{equation}
\mbox{max.}\; z = \sum_{a \in \mathcal{A}} B_a\,x_a
\end{equation}

\paragraph{Modelo completo}\mbox{}\\
\begin{align}
\mbox{max.}\quad & \sum_{a \in \mathcal{A}} B_a\,x_a \notag \\
\mbox{s.a.}\quad 
& \sum_{a \in \mathcal{A}} T_{pa}\,x_a \leq D_p, \quad \forall p \in \mathcal{P} \notag \\
& x_a \geq 0, \quad \forall a \in \mathcal{A}
\end{align}
